% Source: GLM23, arXiv:2203.12563v3, REsubmission.tex, lines 1695--1777
% (block-injective interpolation), with lines 317, 1584--1586, defAgamma,
% Agammasym, and the injective-case fusion argument at lines 1667--1685.

\section{The block-injective interpolation}
\label{sec:mpo_block_mixed_path}

\cite[Section~5]{GarreRubio2023MPOSym} extends the interpolation to states
with a degenerate ground space:
\begin{quote}
    ``We now deal with the case of systems with ground state degeneracy. We
    want to connect two block-injective MPSs (with the same number of
    blocks) that share the $L$-symbols under the action of two MPO algebras
    with the same fusion category $\mathcal{C}$ and same $F$-symbols. For
    this we first have to establish a correspondence between the blocks of
    each MPS, $x:x_0 \to x_1$, and the degeneracies compatible with the
    action of $\mathcal{C}$ to associate
    $\left({L^0}_{abx_0}^{y_0} \right)_{c,k_0 \mu}^{z_0,i_0j_0} =
        \left({L^01}_{abx_1}^{y_1} \right)_{c,k_1 \mu}^{z_1,i_1j_1}$.
    The embedding MPS tensor has the block-diagonal form
    $A(\gamma)=\bigoplus_x A_x(\gamma)$. For every block $x$ the bond
    dimension is $D^0_x+D^1_x$ and it is defined by
    $A_x(\gamma)= A_x \cdot W_x(\gamma)$, [\ldots] where
    $B_x^{pq}= \mathbb{1}_{D^p_x}\otimes \mathbb{1}_{D_x^q}$ and
    $W_x(\gamma) = (1-\gamma)\mathbb{1}_{D_x^0}\oplus
        \gamma\mathbb{1}_{D_x^1}$.''
\end{quote}
The common operators are then described as follows:
\begin{quote}
    ``The MPO tensor $T_a$ is defined as in the injective case, with bond
    dimension $\chi^0_a + \chi^1_a$ for every $a \in \mathcal{C}$, but now we
    have to define the action on all blocks [\ldots] The previous MPO
    tensors define an MPO representation of $\mathcal{C}$ if $L^0=L^1$ and
    $F^0=F^1$, this can be seen as in the injective case. Also, it
    corresponds to an MPO symmetry of $|\psi_{A(\gamma)}\rangle$ since it
    satisfies [\ldots] which is the analogous version of
    Eq.~(fusiontensors2) here, so that the action it is characterized by
    $L^0 = L^1$. The injectivity of every block of $A(\gamma)$ and the
    injectivity of $T_a$ is justified in the same way as in the previous
    case. Moreover, $H(\gamma)$ is well-behaved in the whole path its ground
    space is spanned by $\{ |\psi_{A_x(\gamma)}\rangle\}$.''
\end{quote}

Throughout, the blocks of both endpoints are labelled by one finite set
$x,y,z\in X$, which records the correspondence $x_0\leftrightarrow x_1$,
and the action multiplicities $M_{a,xy}$ are common to both endpoints,
which records the correspondence of the degeneracies. The endpoint blocks
are assumed to occupy separated physical sectors: the local space of the
endpoint $p\in\{0,1\}$ is $\mathcal H_p=\bigoplus_xE^p_x\otimes E^p_x$
with $E^p_x=\C^{D^p_x}$, and the block $A^p_x$ has physical space
$E^p_x\otimes E^p_x$; it is extended by zero on the other summands of
$\mathcal H_p$. This is a restriction beyond the source: line 1584
restricts the local space only to the span of the whole tensor $A_p$, and
for several blocks the separation of that span into block sectors is not
derived here \cite{gap:glm23_block_injective_mixed_path}. Every statement
of this section is made under this restriction. Put $F_x=E^0_x\oplus E^1_x$ and
$\mathcal H=\bigoplus_xF_x\otimes F_x$. The endpoint operators
$T_{p,a}$ act on $\mathcal H_p$ and have bond spaces
$K_{p,a}=\C^{\chi_{p,a}}$; write $T^{yx}_{p,a}$ for the part of $T_{p,a}$
from the summand $x$ to the summand $y$. The action decompositions are
\begin{align}
    (T_{p,a}\cdot A^p_x)^\rho
     & =\sum_{y\in X}\sum_{i=0}^{M_{a,xy}-1}
    W^A_{p,a;xy,i}\,(A^p_y)^\rho\,V^A_{p,a;xy,i},
    \label{eq:mpo_block_action_input}
\end{align}
with maps $V^A_{p,a;xy,i}:K_{p,a}\otimes E^p_x\to E^p_y$ and
$W^A_{p,a;xy,i}:E^p_y\to K_{p,a}\otimes E^p_x$ satisfying
$V^A_{p,a;xy,i}W^A_{p,a;xy',j}=\delta_{(y,i),(y',j)}\Id$. No synthesis
map is assumed to be an adjoint.

\begin{definition}[Block-injective mixed path]
    \label{def:mpo_block_mixed_path}
    \lean{MPSTensor.MPOSymmetry.sectorLetter,
        MPSTensor.MPOSymmetry.sectorTensor,
        MPSTensor.MPOSymmetry.blockMixedPhysicalDim,
        MPSTensor.MPOSymmetry.blockMixedInterpolation,
        MPSTensor.MPOSymmetry.blockMixedTensor}
    \leanok
    \uses{def:mpo_mixed_endpoint_tensor}
    Assume the endpoint blocks occupy separated physical sectors, as
    stated above. For each block $x$, let $A_x(\gamma)$ be the mixed
    interpolation of
    $A^0_x$ and $A^1_x$ of
    Definition~\ref{def:mpo_mixed_endpoint_tensor}, with physical space
    $F_x\otimes F_x$ and bond space $F_x$, extended by zero on the other
    summands of $\mathcal H$. The block-injective mixed tensor is
    $A(\gamma)=\bigoplus_xA_x(\gamma)$, with bond space $\bigoplus_xF_x$.
\end{definition}

\begin{lemma}[Continuity and injectivity of the blocks]
    \label{thm:mpo_block_mixed_injective}
    \lean{MPSTensor.MPOSymmetry.sectorTensor_self,
        MPSTensor.MPOSymmetry.sectorTensor_of_ne,
        MPSTensor.MPOSymmetry.sum_sectorTensor,
        MPSTensor.MPOSymmetry.sum_sigma_sectorTensor,
        MPSTensor.MPOSymmetry.span_range_sectorTensor,
        MPSTensor.MPOSymmetry.isInjective_sectorTensor_iff,
        MPSTensor.MPOSymmetry.sectorTensor_mul_right,
        MPSTensor.MPOSymmetry.continuous_blockMixedInterpolation,
        MPSTensor.MPOSymmetry.isInjective_blockMixedInterpolation}
    \leanok
    \uses{def:mpo_block_mixed_path}
    Each block $A_x(\gamma)$ depends continuously on $\gamma\in\R$. If
    $A^0_x$ and $A^1_x$ are injective, then $A_x(\gamma)$ is injective for
    every $\gamma\in(0,1)$.
\end{lemma}
\begin{proof}
    \leanok
    \uses{thm:mpo_mixed_interior}
    A tensor extended by zero on the other summands has the same span of
    letters as the original tensor. The statements therefore reduce to the
    continuity and interior injectivity of the mixed interpolation of one
    pair of injective tensors.
\end{proof}

\begin{definition}[Common operator of the block-injective path]
    \label{def:mpo_block_mixed_operator}
    \lean{MPSTensor.MPOSymmetry.blockMixedAnalysis,
        MPSTensor.MPOSymmetry.blockMixedSynthesis,
        MPSTensor.MPOSymmetry.blockMixedMPOLetter,
        MPSTensor.MPOSymmetry.blockMixedMPOAnalysis,
        MPSTensor.MPOSymmetry.blockMixedMPOSynthesis,
        MPSTensor.MPOSymmetry.blockMixedMPOSector,
        MPSTensor.MPOSymmetry.operatorSector,
        MPSTensor.MPOSymmetry.blockMixedMPOEntry,
        MPSTensor.MPOSymmetry.blockMixedMPO,
        MPSTensor.MPOSymmetry.blockMixedMPO_finSigmaFinEquiv}
    \leanok
    \uses{def:mpo_mixed_action_tensor,def:mpo_block_mixed_path}
    For each label $a$, define $T_a$ on $\mathcal H$ with bond space
    $K_{0,a}\oplus K_{1,a}$ as follows. The part of $T_a$ from
    $F_x\otimes F_x$ to $F_y\otimes F_y$ has, for $r,k$ in the sector $p$
    of $F_y$ and $F_x$, $s,l$ in the sector $q$, $\alpha\in K_{p,a}$ and
    $\beta\in K_{q,a}$, the entries
    \begin{align}
        (T_a^{(r,s),(k,l)})_{\alpha\beta}
         & =\begin{cases}
                ((T^{yx}_{p,a})^{(r,s),(k,l)})_{\alpha\beta}, & p=q,    \\
                \displaystyle\sum_{i=0}^{M_{a,xy}-1}
                (W^A_{p,a;xy,i})_{(\alpha,k),r}
                (V^A_{q,a;xy,i})_{s,(\beta,l)},               & p\ne q,
            \end{cases}
        \label{eq:mpo_block_mixed_operator}
    \end{align}
    and every entry with unmatched sectors vanishes. The analysis and
    synthesis maps $V_{a;xy,i}$ and $W_{a;xy,i}$ are the direct sums of
    the endpoint maps on matching sectors and vanish on unmatched ones.
\end{definition}

\begin{lemma}[Exact action on every block]
    \label{thm:mpo_block_mixed_action}
    \lean{MPSTensor.MPOSymmetry.blockMixedAnalysis_mul_synthesis,
        MPSTensor.MPOSymmetry.blockMixedMPOAnalysis_mul_synthesis,
        MPSTensor.MPOSymmetry.blockMixedMPOAnalysis_mul_weight,
        MPSTensor.MPOSymmetry.blockMixedMPOLetter_action,
        MPSTensor.MPOSymmetry.blockMixedMPOSector_action,
        MPSTensor.MPOSymmetry.isBiorthogonalDecomposition_blockMixedInterpolation,
        MPSTensor.MPOSymmetry.mpo_mulVec_blockMixedInterpolation}
    \leanok
    \uses{def:mpo_block_mixed_operator,def:glm_boundary_decomposition}
    Assume (\ref{eq:mpo_block_action_input}). Then for every real $\gamma$,
    every label $a$ and every block $x$,
    \begin{align}
        (T_a\cdot A_x(\gamma))^\rho
         & =\sum_{y\in X}\sum_{i=0}^{M_{a,xy}-1}
        W_{a;xy,i}\,A_y(\gamma)^\rho\,V_{a;xy,i},
        \label{eq:mpo_block_mixed_action}
    \end{align}
    and these maps are biorthogonal. Consequently, for every $N\geq1$, the
    periodic operator of $T_a$ maps
    $\ket{V^{(N)}(A_x(\gamma))}$ to
    $\sum_yM_{a,xy}\ket{V^{(N)}(A_y(\gamma))}$.
\end{lemma}
\begin{proof}
    \leanok
    Both sides vanish unless the physical letter lies in some summand
    $F_y\otimes F_y$, and there the sum over the input letters runs over
    the summand $F_x\otimes F_x$ only. On that pair of summands the
    identity is the same entrywise computation as in
    Lemma~\ref{thm:mpo_mixed_action_exact}, with input block $x$ and
    output block $y$: its diagonal sectors are the parts of
    (\ref{eq:mpo_block_action_input}) from $x$ to $y$, and its cross sectors
    are the defining contractions (\ref{eq:mpo_block_mixed_operator}). The
    analysis maps commute with the weights $W_x(\gamma)$ and $W_y(\gamma)$
    as in that lemma. Taking the trace of a word, biorthogonality gives the
    periodic statement.
\end{proof}

For the fusion, fix a finite label set with fusion multiplicities
$N_{ab}^c$, and fusion decompositions
\begin{align}
    (T_{p,a}T_{p,b})^{\rho\sigma}
     & =\sum_{c}\sum_{\mu=0}^{N_{ab}^c-1}
    W^F_{p,ab;c\mu}T_{p,c}^{\rho\sigma}V^F_{p,ab;c\mu}
    \label{eq:mpo_block_fusion_input}
\end{align}
with biorthogonal maps, as in Section~\ref{sec:mpo_mixed_fusion}. For
blocks $x,y$, let $L_{p,ab;xy}$ be the comparison matrix of the
sequential action paths $x\to z\to y$ with the fusion-then-action paths
$x\to y$, defined from the maps of (\ref{eq:mpo_block_action_input}) and
(\ref{eq:mpo_block_fusion_input}) as in
Definition~\ref{def:glm_boundary_action_l_matrix}.

\begin{lemma}[Block comparison relations]
    \label{thm:mpo_block_mixed_fusion_cross}
    \lean{MPSTensor.MPOSymmetry.wordTupleSpanTop_sectorTensor,
        MPSTensor.MPOSymmetry.blockSequentialAnalysis,
        MPSTensor.MPOSymmetry.blockSequentialSynthesis,
        MPSTensor.MPOSymmetry.blockFusionThenActionAnalysis,
        MPSTensor.MPOSymmetry.blockFusionThenActionSynthesis,
        MPSTensor.MPOSymmetry.blockSequentialSynthesis_apply,
        MPSTensor.MPOSymmetry.blockSequentialAnalysis_apply,
        MPSTensor.MPOSymmetry.blockFusionThenActionSynthesis_apply,
        MPSTensor.MPOSymmetry.blockFusionThenActionAnalysis_apply,
        MPSTensor.MPOSymmetry.block_actionLMatrix_relations,
        MPSTensor.MPOSymmetry.blockCrossEndpoint_actionTree_contraction,
        MPSTensor.MPOSymmetry.blockCrossEndpoint_actionTree_contraction_apply,
        MPSTensor.MPOSymmetry.blockCrossEndpoint_fusion_entry}
    \leanok
    \uses{thm:mpo_mixed_fusion_cross,thm:glm_boundary_action_l_exists,
        def:glm_boundary_action_l_matrix}
    Assume that all blocks $A^p_x$ are injective with $D^p_x>0$, and
    that $L_{0,ab;xy}=L_{1,ab;xy}$ for a pair of labels $a,b$ and a pair
    of blocks $x,y$. Then the crossed contraction
    (\ref{eq:mpo_mix_fusion_cross}) holds for the sequential and
    fusion-then-action maps from $x$ to $y$, for every rectangular matrix
    $X:E^1_y\to E^0_y$, with the intermediate block $z$ summed on the
    sequential side.
\end{lemma}
\begin{proof}
    \leanok
    Since every letter of $\bigoplus_x A^p_x$ lies in one summand of
    $\mathcal H_p$ and every block is injective, the letters of length one
    already span $\bigoplus_x M_{D^p_x}$. The comparison relations of
    Lemma~\ref{thm:mpo_mixed_fusion_cross} hold therefore block by block,
    and the crossed contraction follows from them by the same computation.
\end{proof}

\begin{lemma}[Fusion of the common operators]
    \label{thm:mpo_block_mixed_fusion}
    \lean{MPSTensor.MPOSymmetry.operatorSector_fusion_entry,
        MPSTensor.MPOSymmetry.blockMixedMPOLetter_fusion,
        MPSTensor.MPOSymmetry.isBiorthogonalDecomposition_blockMixedMPO,
        MPSTensor.MPOSymmetry.isMPOFusionAlgebra_blockMixedMPO}
    \leanok
    \uses{thm:mpo_block_mixed_fusion_cross,thm:mpo_mixed_fusion_maps,
        def:mpo_mixed_fusion_maps,def:mpo_fusion_algebra}
    Assume (\ref{eq:mpo_block_action_input}),
    (\ref{eq:mpo_block_fusion_input}), that every block $A^p_x$ is
    injective with $D^p_x>0$, and that $L_{0,ab;xy}=L_{1,ab;xy}$ for a pair
    of labels $a,b$ and all blocks $x,y$. Then
    \begin{align}
        (T_aT_b)^{\rho\sigma}
         & =\sum_c\sum_{\mu=0}^{N_{ab}^c-1}
        W^F_{ab;c\mu}T_c^{\rho\sigma}V^F_{ab;c\mu},
        \label{eq:mpo_block_mixed_fusion}
    \end{align}
    with the biorthogonal direct sums $V^F_{ab;c\mu}$ and $W^F_{ab;c\mu}$
    of the endpoint fusion maps. If the comparison matrices agree for all
    labels and blocks, the periodic operators of the family $T_a$ satisfy
    the fusion rules $N_{ab}^c$ at every length.
\end{lemma}
\begin{proof}
    \leanok
    The product $T_aT_b$ from $F_x\otimes F_x$ to $F_y\otimes F_y$ sums over
    the intermediate summands $F_z\otimes F_z$. On the diagonal sectors
    the identity is the part of (\ref{eq:mpo_block_fusion_input}) from $x$
    to $y$, with the intermediate letter split by its summand. On the cross
    sectors it is the crossed contraction of
    Lemma~\ref{thm:mpo_block_mixed_fusion_cross} evaluated at a matrix unit,
    exactly as in Lemma~\ref{thm:mpo_mixed_fusion_exact}. All other
    sectors vanish on both sides. The fusion rules follow by taking traces
    of words.
\end{proof}

The common $F$-symbols are not needed for the exact fusion
(\ref{eq:mpo_block_mixed_fusion}).

The path Hamiltonian is the two-site canonical parent Hamiltonian
$H(\gamma)$ of $A(\gamma)$ on a ring of $N$ sites, as in the injective
case.

\begin{lemma}[Ground space of the block-injective path]
    \label{thm:mpo_block_mixed_ground_space}
    \lean{MPSTensor.MPOSymmetry.evalWord_sectorTensor_map,
        MPSTensor.MPOSymmetry.mpv_sectorTensor_self,
        MPSTensor.MPOSymmetry.mpv_sectorTensor_of_ne,
        MPSTensor.MPOSymmetry.blockPeriodicMpvMapES_sectorTensor_injective,
        MPSTensor.MPOSymmetry.wordTupleSpanTop_blockMixedInterpolation,
        MPSTensor.MPOSymmetry.blockPeriodicMpvMapES_blockMixedInterpolation_injective,
        MPSTensor.MPOSymmetry.ker_hamiltonian_parentRingInteraction_blockMixedTensor,
        MPSTensor.MPOSymmetry.linearIndependent_mpv_blockMixedInterpolation,
        MPSTensor.MPOSymmetry.continuous_blockMixed_groundSpace_starProjection,
        MPSTensor.MPOSymmetry.exists_uniform_blockMixed_hamiltonian_parentRingInteraction_gap,
        MPSTensor.MPOSymmetry.tendsto_blockMixed_groundSpace_starProjection}
    \leanok
    \uses{def:mpo_block_mixed_path,def:parent_hamiltonian_es}
    Assume the endpoint blocks occupy separated physical sectors, that there
    is at least one block, and that every block $A^p_x$ is injective with
    $D^p_x>0$. Let $N\geq2$.
    \begin{enumerate}
        \item For $\gamma\in(0,1)$ the ground space of $H(\gamma)$ is
              spanned by the block states $\ket{V^{(N)}(A_x(\gamma))}$.
        \item For $\gamma\in[0,1]$ the block states are linearly
              independent.
        \item The ground projection of $H(\gamma)$ is continuous on
              $(0,1)$, and every compact subset of $(0,1)$ admits one
              $\delta>0$ such that, for all $N\geq2$, every nonzero
              eigenvalue of $H(\gamma)$ is at least $\delta$.
        \item As $\gamma\to\gamma_0\in[0,1]$ within $(0,1)$, the ground
              projection of $H(\gamma)$ converges to the orthogonal
              projection onto the span of the block states at $\gamma_0$.
    \end{enumerate}
\end{lemma}
\begin{proof}
    \leanok
    \uses{thm:mpo_block_mixed_injective,thm:mpo_block_mixed_fusion_cross,
        thm:mpo_mixed_interior,thm:mpo_mixed_periodic_line_continuity,
        thm:block_ground_space_continuity,thm:compact_block_parent_gap,
        lem:uncle_local_projector_continuity}
    For $\gamma\in(0,1)$ every block is injective, and since the blocks
    occupy separated physical sectors, the letters of length one span
    $\bigoplus_xM_{D^0_x+D^1_x}$. The block-injective parent Hamiltonian
    theory with simultaneous injectivity length one and a two-site
    interaction identifies the ground space with the span of the block
    states, gives continuity of its projection, and gives the uniform gap
    on compact parameter sets. For $N\geq1$, a word $\tau$ of length $N$ in
    the basis of $F_y\otimes F_y$, and its image $\iota_y(\tau)$ in the basis
    of $\mathcal H$,
    \begin{align}
        V^{(N)}(A_x(\gamma))_{\iota_y(\tau)}
         & =\delta_{xy}\,V^{(N)}\bigl(A_y(\gamma)|_{F_y\otimes F_y}\bigr)_\tau,
        \label{eq:mpo_block_mixed_sector_support}
    \end{align}
    since for $y\neq x$ the first letter of $\iota_y(\tau)$ is a zero matrix
    of $A_x(\gamma)$. The states
    $V^{(N)}(A_y(\gamma)|_{F_y\otimes F_y})$ are nonzero for
    $\gamma\in[0,1]$, by the injective case. Evaluating
    $\sum_xc_xV^{(N)}(A_x(\gamma))=0$ at $\iota_y(\tau)$ through
    (\ref{eq:mpo_block_mixed_sector_support}) gives
    $c_yV^{(N)}(A_y(\gamma)|_{F_y\otimes F_y})_\tau=0$ for every $\tau$, so
    $c_y=0$: the block states are linearly independent. The span of a
    continuous family of linearly independent vectors has a projection that
    is continuous at every such point, and it
    agrees with the ground projection on $(0,1)$.
\end{proof}

The claim that the operators $T_a$ are injective is not treated in this
section, nor are the kernels and gaps of the endpoint Hamiltonians. The
source does not show that $H(\gamma)$ commutes with the periodic operators
of $T_a$, and the one-block case of this construction is the injective
interpolation, for which the canonical parent need not commute with a
non-unitary symmetry; see \cite{gap:glm23_block_injective_mixed_path}.
